\section{Actual bounds for the retained four-window costs}

Continue with (HT1)--(HT18), now with $k\geq3$.  All unadorned
source norms and quotient volumes in this section are those of the original
$m_{h,k}(u)\,du$, with $g=2\xi$ and the complete hypothetical quartet.
Write $\mathcal L=\mathcal L_{h,k}$ and retain $q=\ell_k(k+1)^2$.
The threshold-four lower bound is inherited from
\emph{Consecutive identities and the first window}, (CJ.9)--(CJ.19);
it is already explicitly present in the recovered transcript at lines
4202--4228.  It is not a new claim of the present continuation.

\subsection{Every term of the inherited lower bound}

Put $d_j=V_j/V_{j-1}$ for $j\geq q$.  For the two original windows
$I_-=(q-1,2q-1)$ and $I_+=(q,2q)$, retain the contraction penalties
\[
 \begin{split}
 P_-&=-\log(1-d_{2q-1})-2\sum_{j=q}^{2q-2}\log(1-d_j),\\
 P_+&=-\log(1-d_q)-\log(1-d_{2q})
       -2\sum_{j=q+1}^{2q-1}\log(1-d_j).
 \end{split}\tag{HC1}
\]
The source identities prove $0<d_j<1$ in these ranges.  The real phase
$\phi_N$ is defined by the retained source pairing
$b_N^*G_Nb_{N+1}=i\omega_N\phi_N$.  The first window uses its original
boundary identity; no inverse at degree $q-2$ occurs.  Define weights
\[
 w_{q-1}=w_{2q-1}=1,\qquad w_N=2\quad(q\leq N\leq2q-2),
 \qquad\sum_{N=q-1}^{2q-1}w_N=2q.
\]
Adding the two exact source products gives
\[
 \mathcal B_k:=\log\frac{V_{q-1}V_q}{V_{2q-1}V_{2q}}
 =\sum_{N=q-1}^{2q-1}w_N\log(\varepsilon_N^2+\phi_N^2)
 -\log\frac{\omega_{2q-1}\omega_{2q}}{\omega_{q-1}\omega_q}
 +P_-+P_+.
 \tag{HC2}
\]
Here every $\varepsilon_N\geq\mathcal L$, including the first degree.

For clarity, retain all constants in the inherited comparable-window
bound.  With $\alpha=\pi/2$, $B=42+8m$, and
$\vartheta_h=\int_{-1}^{1}w_h$, put
\[
 \begin{gathered}
 t_h=\min(1,\vartheta_h),\quad c_0=\min(1,2c_h/3),\quad
 \ell_h=\tfrac12 e^{-\alpha}\sqrt{c_0t_h}\,2^{-B/2},\\
 X=\max(1,M_h(b),M_h(-b)),\quad U_h=4\sqrt X/b,\\
 C_h^{\rm win}=27\max(1,U_h)^3\max(1,\ell_h^{-1})^2,
 \qquad D_h=\frac{\delta}{2C_h^{\rm win}}.
 \end{gathered}\tag{HC3}
\]
Here $0<b<\pi/2$, and $c_h>0$ is the original proved convolution
constant.  The source proves
$(\omega_{n+r}/\omega_n)^{1/(2r)}\leq C_h^{\rm win}n$ for
$n\geq k\geq3$, $n/2\leq r\leq2n$.  Thus it applies to both literal
windows, and $0<D_h<1/108$.

The exact residual in (HC2) has the following five nonnegative pieces:
\[
 \begin{split}
 \mathcal R_k={}&P_-+P_+
 +\sum_Nw_N\log(1+\phi_N^2/\mathcal L^2)\\
 &+\sum_Nw_N\log\frac{\varepsilon_N^2+\phi_N^2}
                              {\mathcal L^2+\phi_N^2}\\
 &+4q\log(C_h^{\rm win}q)
       -\log\frac{\omega_{2q-1}\omega_{2q}}{\omega_{q-1}\omega_q}
 +4q\log\frac{2\mathcal L}{\delta kq}.
 \end{split}\tag{HC4}
\]
Both sums use exactly $q-1\leq N\leq2q-1$.  Nonnegativity follows,
respectively, from $0<d_j<1$, the squared phases, (HT9), the two
comparable-window estimates, and $\mathcal L/q\geq\delta k/2$.
Expanding each logarithm in (HC4), with $D_h=\delta/(2C_h^{\rm win})$,
proves the exact identity
\[
 \boxed{\mathcal B_k=4q\log(D_hk)+\mathcal R_k,
 \qquad\mathcal R_k\geq0.}\tag{HC5}
\]
This retains, rather than discards, all the contraction, phase, norm,
and trace savings in the original lower bound.

\subsection{A proved upper bound for these same four original volumes}

The original convolution estimate and Laplace moments are
\[
 m_{h,k}(u)\geq c_h\vartheta_h^{k-3}
 e^{-\alpha(|u|+k-3)}(1+|u|+k-3)^{-B},\qquad
 M_h(\pm b)=\int_{\mathbb R}e^{\pm bu}w_h(u)\,du.
\]
Write $\mu_h=\int_{\mathbb R}w_h(u)\,du$ for the original source mass.
Put
\[
 \begin{gathered}
 R_0=\sqrt{\delta^2+\gamma^2},\quad
 D=\max(1,2/b,2R_0),\quad T=Dq,\quad r=kR_0/T\leq1/2,\\
 a_k(T)=c_h\vartheta_h^{k-3}e^{-\alpha(T+k-3)}
 (1+T+k-3)^{-B},\\
 M_k=M_h(b)^k+M_h(-b)^k,\qquad
 U_k(T)=\mu_h^k+\frac{(2q-2)!}{(bT)^{2q-2}}M_k,\\
 A_N=\sqrt{\frac{2N+1}{2}}\frac{4^{N+1}-1}{3},\qquad
 B_q(r)=\left(\frac{1+r}{1-r}\right)^q,\qquad
 C_N=\frac{U_k(T)A_N^2B_q(r)^2}{T a_k(T)}.
 \end{gathered}\tag{HC6}
\]
The capital $D$ in the interval scale is distinct from the inherited
lower constant $D_h$ and from the source differential.  All masses,
coordinates, and root multiplicities remain as in (HT1)--(HT7).
One has the actual finite Hermitian comparisons
\[
 G_{q-1}\preceq C_{2q-1}G_{2q-1},\qquad
 G_q\preceq G_{q-1}\preceq C_{2q}G_{2q}.
\]
Consequently, with the exact ratio
\[
 \rho_q=\frac{4q+1}{4q-1}
 \left(\frac{4^{2q+1}-1}{4^{2q}-1}\right)^2,
\]
\[
 \boxed{0\leq\mathcal B_k
 \leq q\log C_{2q-1}+q\log C_{2q}
 =2q\log C_{2q-1}+q\log\rho_q.}\tag{HC7}
\]

\paragraph{Proof of the degree extension and its maps.}
Use the exact substitution $f(x)=P(c+iTx)$ with inverse
$P(S)=f((S-c)/(iT))$.  Its source measure is $T m_{h,k}(Tx)\,dx$,
with mass $\mu_h^k$, and its monic relation is
$\chi_T(x)=(iT)^{-q}\chi(c+iTx)$.  The roots $z_1,\ldots,z_q$,
listed with their repetitions, have modulus at most $r$.
The raw jet of order $d$ acquires exactly $(iT)^d$; the quotient
determinant acquires $T^{-q(q-1)}$.  This common factor cancels in
each displayed quotient-volume ratio.

The Legendre expansion on $[-1,1]$ gives
$\|f\|_{{\rm coeff},1}\leq A_N\|f\|_{L^2([-1,1])}$ for
$\deg f\leq N$.  Indeed Rodrigues' polynomial has coefficient norm
at most $\binom{2j}{j}\leq4^j$ and squared integral $2/(2j+1)$.
Pairing the expansion with this polynomial bounds its coefficient by
$\sqrt{(2j+1)/2}\|f\|_2$; summing the geometric series gives exactly
$A_N$ in (HC6).

Write $\prod_a(1-z_at)^{-1}=\sum_{j\geq0}h_jt^j$.
Geometric-series multiplication gives
$|h_j|\leq\binom{q+j-1}{j}r^j$ and
$\|\chi_T\|_{{\rm coeff},1}\leq(1+r)^q$.
For every integer $s\geq0$, monic division of $x^{q+s}$ has quotient
$\sum_{j=0}^sh_jx^{s-j}$.  The inverse-series identity cancels all
other terms of degree at least $q$.  Removing the leading coefficient
one therefore gives
\[
 \|\operatorname{rem}_{\chi_T}x^{q+s}\|_{{\rm coeff},1}
 \leq(1+r)^q\sum_{j=0}^s\binom{q+j-1}{j}r^j-1
 \leq B_q(r).
\]
The bound is valid also for $s=q$, the extra degree required here.
At degrees below $q$ the remainder is the original monomial, whose
coefficient norm is one.  Linearity gives the same remainder-map
bound for every degree at most $2q$.

For a remainder $v$ of degree below $q$,
\[
 \int |v(u/T)|^2m_{h,k}(u)\,du
 \leq U_k(T)\|v\|_{{\rm coeff},1}^{2}.
\]
This follows by bounding each monomial by
$\max(1,|u/T|^{q-1})$ and using
$\int u^{2q-2}m_{h,k}\leq(2q-2)!b^{-(2q-2)}M_k$.
Restriction of the original norm to $[-T,T]$ gives the lower bound
$T a_k(T)\|f\|_{L^2([-1,1])}^2$.  Combining these three estimates proves
\[
 \|\operatorname{rem}_{\chi}P\|_{H_{q-1}}^2
 \leq C_N\|P\|_{H_N}^2\qquad(N=2q-1,2q).
\]
The map is the literal remainder representative of the original jet.
Apply it to the canonical least lift of $x\in E$ at degree $N$.
Its norm at degree $q-1$ is at least $x^*G_{q-1}x$; this proves the
Hermitian comparisons preceding (HC7).  Inclusion of polynomial degrees
gives $G_q\preceq G_{q-1}$ and monotonicity of all volumes.  Taking
determinants proves (HC7).  Applying the same remainder inequality to
any nonzero polynomial below degree $q$ proves $C_N\geq1$ without an
additional assumption.\hfill$\square$

For an entirely scalar upper bound retain the original constants
\[
 K=\max(1,\vartheta_h^{-1}),\qquad
 C_{\rm env}=\max(1,2/(3c_hD)),\quad a=\alpha D+\log256,
\]
and the original expression (AU.31)
\[
 \mathcal U_k=a q^2+
 \left(\log(XK)+\alpha+\frac{16R_0}{3D}\right)kq
 +Bq\log(1+(D+1)q)+q\log C_{\rm env}.
\]
Then the proved four-window upper bound is
\[
 \boxed{\mathcal B_k\leq\widehat{\mathcal U}_k
 :=2\mathcal U_k+q\log20.}\tag{HC8}
\]
To verify the constants, $bT\geq2q$ implies the factorial ratio in
(HC6) is at most one; also $\mu_h\leq X$ and $M_k\leq2X^k$, so
$U_k(T)\leq3X^k$.  Directly from (HC6),
\[
 A_{2q-1}^2\leq(2q/9)256^q,\qquad
 A_{2q}^2\leq(40q/9)256^q.
\]
The logarithmic derivative $2/(1-r^2)\leq8/3$ on $[0,1/2]$ gives
$\log B_q(r)^2\leq16kR_0/(3D)$.  Substitution into $C_N$ yields
$q\log C_{2q-1}\leq\mathcal U_k$ and
$q\log C_{2q}\leq\mathcal U_k+q\log20$.
Adding proves (HC8).  This argument does not use the false inequality
$\rho_2<20$; in the actual quartet domain $q\geq16$ anyway.

\subsection{The exact compatible interval for the signed arithmetic correction}

Let $V_N^\Gamma$ denote the inherited Gamma-reference quotient volume
for the \emph{same} original polynomial $\chi$, and define
\[
 T_N=V_N/V_N^\Gamma,\qquad
 \mathcal B_k^\Gamma=\log\frac{V_{q-1}^\Gamma V_q^\Gamma}
 {V_{2q-1}^\Gamma V_{2q}^\Gamma},\qquad
 \Delta_k^\Gamma=\log\frac{T_{q-1}T_q}{T_{2q-1}T_{2q}}.
\]
Cancellation of the four literal volume factors gives
$\mathcal B_k=\mathcal B_k^\Gamma+\Delta_k^\Gamma$.
No auxiliary polynomial is introduced and no sign is imposed on
$\Delta_k^\Gamma$. The exact coordinate identification with (ACM1)--(ACM5)
is $y=u$, with $S=k/2+iu$ and unchanged ascending coefficients.
Its arithmetic form is the original $m_{h,k}=w_h^{*k}$, while the
reference form is precisely (ACM2), including its full convolution
mass. The relation matrices are the same multiplication by the full
$\chi$ of (HT2). Thus apply (ACM11a)--(ACM16) to this pair and write
$\beta_{p,\Gamma}^\pm$ for the resulting simultaneous endpoints,
whose initial value is $F_0=\mathcal B_k^\Gamma$. In particular
$\beta_{p,\Gamma}^-\leq\mathcal B_k\leq\beta_{p,\Gamma}^+$.
This uses every spectral, trace-norm, overlap, variance, angle and
constructed signed-residual control on the original Gamma-to-arithmetic
path. Put $T_k=4q\log(D_hk)$ for this scalar expression only. Then
(HC5), (HC8) and these endpoints give the current intervals
\[
 \boxed{
 \max\{0,-T_k,\beta_{p,\Gamma}^--T_k\}\leq\mathcal R_k
 \leq\min\{\widehat{\mathcal U}_k,\beta_{p,\Gamma}^+\}-T_k.}
 \tag{HC9}
\]
\[
\begin{split}
 \max\{0,T_k,\beta_{p,\Gamma}^-\}-\mathcal B_k^\Gamma
 &\leq\Delta_k^\Gamma\\
 &\leq\min\{\widehat{\mathcal U}_k,\beta_{p,\Gamma}^+\}
                          -\mathcal B_k^\Gamma.
\end{split}
\tag{HC10}
\]
Every entry follows by subtracting the same displayed initial value
from a proved bound on $\mathcal B_k$. Removing the source-dependent
$\beta_{p,\Gamma}^\pm$ entries gives the earlier coarse intervals.
Those coarse intervals have strictly positive width for every $k\geq3$.

Indeed all nonquadratic terms of $\mathcal U_k$ are nonnegative,
$a>1$, $q\geq(k+1)^2>2k$, and $D_h<1$.  Thus
\[
 \widehat{\mathcal U}_k\geq2a q^2
 >4qk\geq\max\{0,4q\log(D_hk)\}.
\]
Moreover, keeping the fixed packet constants, each lower-order term
divided by $q^2$ tends to zero, so
\[
 \lim_{k\to\infty}
 \frac{\widehat{\mathcal U}_k-\max\{0,4q\log(D_hk)\}}{q^2}
 =2(\alpha D+\log256)>0.\tag{HC11}
\]
This proves the stated width of the coarse bound functions obtained
by removing the source-dependent entries in (HC9)--(HC10). The
refined intersection retains those computed entries; (HC11) makes
no assertion that its width equals the coarse width.  It does not construct a packet realizing an arbitrarily
chosen point of either interval, and it does not assert the existence
of the hypothetical zero.  In particular the interval may lie entirely
at negative values of $\Delta_k^\Gamma$; determinant-ratio positivity
does not imply positivity of its signed four-window logarithm.

\subsection{The holonomy energy is tied to the same endpoint contractions}

Fix an actual holonomy $\theta$ and period $L$.  In this paragraph all
metrics, lifts and volumes carry that same phase.  Put
\[
 \begin{gathered}
 e_N=\log(V_N/V_{N+1}),\quad
 \beta_N=\lambda_{\max}
 (G_{N+1}^{-1/2}G_NG_{N+1}^{-1/2}),\\
 J_N^{D}=\|D_{L,\theta}R_NG_N^{-1/2}\|_{\rm HS}^{2},\quad
 \mathcal E_N=\|\mathscr B_NG_N^{-1/2}\|_{\rm HS}^{2}.
 \end{gathered}
\]
These use original source moments through degree $2N+2$.  They obey
\[
 \boxed{
 \frac{\mathcal L^2}{2}\leq\mathcal E_N
 \leq(1+\sqrt{\beta_N})^2J_N^D,
 \qquad
 e_N\geq2\log\max\left\{1,
 \frac{\mathcal L}{\sqrt{2J_N^D}}-1\right\}.}
 \tag{HC12}
\]
The denominator is positive: integration by parts on the stated
twisted domain gives
$\|D_{L,\theta}f\|^2=c^2\|f\|^2+\|\partial_rf\|^2$, whence
$J_N^D\geq c^2q>0$.

To prove the upper inequality, $D_{L,\theta}R_Nv$ is an original
degree-$(N+1)$ representative of $Av$.  The same quotient minimum
therefore gives
\[
 A^*G_{N+1}A\preceq R_N^*D_{L,\theta}^*D_{L,\theta}R_N.
\]
Use $G_N\preceq\beta_NG_{N+1}$, sandwich by $G_N^{-1/2}$, and take
traces.  This gives
$\|R_NAG_N^{-1/2}\|_{\rm HS}^2\leq\beta_NJ_N^D$.
The triangle inequality in $\mathscr B_N=D_{L,\theta}R_N-R_NA$
proves the upper inequality in (HC12); (HT14) proves the lower.
All eigenvalues in the definition of $\beta_N$ are at least one by
degree monotonicity.  In fact the original rank-one kernel update proves
the exact identity
\[
 e_N=\log\beta_N,\qquad
 \beta_N=1+\frac{b_{N+1}^*G_Nb_{N+1}}{\omega_{N+1}}.
 \tag{HC12a}
\]
Here $b_{N+1}$ and $\omega_{N+1}$ are the remainder column and monic
squared norm for this same original phase measure.  Indeed
$K_{N+1}=K_N+b_{N+1}b_{N+1}^*/\omega_{N+1}$, so
\[
 K_N^{-1/2}K_{N+1}K_N^{-1/2}=I+vv^*,\qquad
 v=K_N^{-1/2}b_{N+1}/\sqrt{\omega_{N+1}}.
\]
The relative matrix defining $\beta_N$ is
$K_{N+1}^{1/2}K_N^{-1}K_{N+1}^{1/2}=X^*X$ for the invertible
$X=K_N^{-1/2}K_{N+1}^{1/2}$.  The displayed matrix $I+vv^*$ is
$XX^*$, and $X(X^*X)X^{-1}=XX^*$ proves they have identical eigenvalues.
These are $1$ with multiplicity at least $q-1$ and $1+\|v\|^2$.
Taking the determinant gives
$V_N/V_{N+1}=\det K_{N+1}/\det K_N=1+\|v\|^2=\beta_N$,
which proves (HC12a), including a zero update if one occurs.
Rearranging the two energy bounds now proves the second inequality
of (HC12) with this exact endpoint identity.

The same weights as in (HC2) reconstruct the actual phase volume:
\[
 \boxed{
 \mathcal B_{k,\theta}=\sum_{N=q-1}^{2q-1}w_N e_N
 \geq2\sum_{N=q-1}^{2q-1}w_N
 \log\max\left\{1,\frac{\mathcal L}{\sqrt{2J_N^D}}-1\right\}.}
 \tag{HC13}
\]
Thus a small original derivative-source energy at a degree has the
explicit endpoint contraction cost (HC12); that cost enters the same
four-window budget, with its actual multiplicity.

Finally the supplied holonomy construction allows any fixed
$0<\eta<1$ by choosing the original period $L$ from its computed
weighted source comparison at degree $2q+2$. We now apply that
comparison on the actual source before quotient minimization.
Let $I_{2q}^{2q+2}:\mathcal P_{2q}\hookrightarrow\mathcal P_{2q+2}$
be the literal coefficient inclusion and retain $H=\mathcal P_{2q}$.
The exact forms on $H$ are
\[
\begin{split}
 M_0^\theta&=(I_{2q}^{2q+2})^*M_{2q+2}I_{2q}^{2q+2}
       =\int v(u)^*v(u)m_{h,k}(u)\,du,\\
 M_1^\theta&=(I_{2q}^{2q+2})^*M_{2q+2,\theta}I_{2q}^{2q+2}
       =\frac{2\pi}{L}\sum_{r\in\mathbb Z}
          m_{h,k}(u_r)v(u_r)^*v(u_r),\\
 v(u)&=(1,c+iu,\ldots,(c+iu)^{2q}),\quad c=k/2,
       \qquad u_r=(2\pi r+\theta)/L.
\end{split}
\tag{HC14a}
\]
This is (HT6) with its original mass and lattice. In terms of the
actual source maps the two forms are the squared norms of
$\mathcal U_k\mathcal V|_H$ and
$\mathcal Z_{L,\theta}\mathcal U_k\mathcal V|_H$, respectively.
Thus no relative source variable, full Taylor unit or coefficient
is removed. Restricting the proved (PSC.33) source inequality by
$I_{2q}^{2q+2}$ proves
$(1-\eta)M_0^\theta\preceq M_1^\theta\preceq(1+\eta)M_0^\theta$.
For each polynomial degree $q-1\leq N\leq2q$, both quotient minima are over
the identical fibre $J_NP=v$; the sandwich also gives the earlier
quotient comparisons by taking infima on those fibres.

Set $M_x^\theta=(1-x)M_0^\theta+xM_1^\theta$ for $x\in[0,1]$.
Use the literal $I_N:\mathcal P_N\to H$ and the unpadded relation
$B_N^{\rm AW}:\mathcal P_{N-q}\to\mathcal P_N$, $P\mapsto\chi P$,
in (ACM4), with $M_x=M_x^\theta$. Thus the full relation map is
$B_N^\theta=I_NB_N^{\rm AW}:\mathcal P_{N-q}\to H$ and the
resulting projections are exactly
\[
\begin{split}
 P_N^\theta&=I_N(I_N^*M_x^\theta I_N)^{-1}I_N^*M_x^\theta,\\
 Q_N^\theta&=B_N^\theta((B_N^\theta)^*M_x^\theta B_N^\theta)^{-1}
                                      (B_N^\theta)^*M_x^\theta.
\end{split}
\]
The negative relation domain is the zero space. Put
\[
\begin{split}
 U_x^\theta&=Q_{2q-1}^\theta+Q_{2q}^\theta-Q_q^\theta,\\
 W_x^\theta&=P_{2q-1}^\theta-P_{q-1}^\theta
                         +P_{2q}^\theta-P_q^\theta,\\
 C_x^\theta&=(M_x^\theta)^{-1}(M_1^\theta-M_0^\theta),\qquad
 D_\theta=(M_0^\theta)^{-1}M_1^\theta.
\end{split}
\tag{HC14b}
\]
Differentiating each canonical quotient Gram exactly as in (ISM18)
gives
$F_\theta'(x)=\operatorname{Tr}(C_x^\theta(U_x^\theta-W_x^\theta))$,
where $F_\theta(0)=\mathcal B_k$ and
$F_\theta(1)=\mathcal B_{k,\theta}$. In particular the zero relation
domain still gives $Q_{q-1}^\theta=0$. The two complete spectra are
$0^{[q]},1^{[2]},2^{[q-1]}$. The original monomial and relation flags,
and all minimum-section angle entries including zeroes, are exactly
the ones used in the proof of (ACM11a).

Define $J_\theta,I_\theta,j_{p,\theta},e_{p,\theta}$ by the literal
substitution
\[
 (M_x,C_x,U_x,W_x,D,F(x))\longmapsto
 (M_x^\theta,C_x^\theta,U_x^\theta,W_x^\theta,D_\theta,F_\theta(x))
\]
throughout (ACM11a) and (ACM15). The index $p\geq0$ is the finite-series degree;
the original period remains $L$. These are separately calculated
source-phase quantities. Explicitly, $J_\theta$ is the integral of
the pointwise minimum of the full ordered-spectrum expression,
the exact midpoint trace norm, the actual-overlap majorant, the
centered Hilbert--Schmidt expression, and the full-angle expression;
$I_\theta$ divides these bounds by
$(2q-1)\sqrt{1-e^{-F_\theta(x)/(2q-1)}}$ as in (ACM11a).
The original continuous measure and the positive lattice measure
both have positive Grams, as proved in (HT6). Therefore the
$\chi^{\#_k}$ argument in (ACM11a) proves $F_\theta(x)>0$ on this
entire path, including its atomic endpoint. All substitutions in
the derivative and nonlinear calculations are consequently valid.
The signed construction retains the exact finite residual and gives
\[
\begin{gathered}
 j_{p,\theta}-e_{p,\theta}
 \leq\mathcal B_{k,\theta}-\mathcal B_k
 \leq j_{p,\theta}+e_{p,\theta},\\
 0\leq e_{p,\theta}\leq
 \left(\frac{b_n^\theta-b_1^\theta}{b_n^\theta+b_1^\theta}\right)^{p+1}
       \sqrt{K_{D_\theta}(8q-6)}\longrightarrow0,
 \qquad n=2q+1.
\end{gathered}
\tag{HC14c}
\]
Here every $b_j^\theta$ is an eigenvalue of $D_\theta$ in $M_0^\theta$,
and $K_{D_\theta}=\sum_{j=1}^n
(|b_j^\theta-1|/\min\{1,b_j^\theta\})^2$. The finite geometric
identity and the original Hilbert--Schmidt trace proof in (ACM15)
prove every term of (HC14c) after the explicit substitution (HC14b).
The source sandwich places each $b_j^\theta$ in $[1-\eta,1+\eta]$.
The simultaneous original phase error is therefore
\[
\begin{split}
 |\mathcal B_{k,\theta}-\mathcal B_k|
 &\leq J_\theta
 \leq\mathcal L_q(D_\theta)
 \leq(2q-1)\log\frac{1+\eta}{1-\eta},\\
 |\mathscr H_{2q-1}(\mathcal B_{k,\theta})
       -\mathscr H_{2q-1}(\mathcal B_k)|
 &\leq I_\theta\leq\log\frac{1+\eta}{1-\eta}.
\end{split}
\tag{HC14}
\]
Thus the earlier $2q$ determinant-only allowance follows as a
larger bound; every stronger pointwise control and signed center
has been retained before any scalar allowance is taken.

For any integers $p,r\geq0$, put
\[
\begin{split}
 U_{k,p}^{\rm ar}&=\min\{\widehat{\mathcal U}_k,
                                       \beta_{p,\Gamma}^+\},\\
 U_{k,\theta;p,r}&=\min\bigl\{
 U_{k,p}^{\rm ar}+J_\theta,
 \Psi_{2q-1}(\mathscr H_{2q-1}(U_{k,p}^{\rm ar})+I_\theta),
 U_{k,p}^{\rm ar}+j_{r,\theta}+e_{r,\theta}\bigr\}.
\end{split}
\tag{HC14d}
\]
The actual $\mathcal B_k$ is positive and at most $U_{k,p}^{\rm ar}$
by (HC10). The additive and signed expressions in (HC14d) increase
with their initial value. Both $\mathscr H_{2q-1}$ and
$\Psi_{2q-1}$ increase on the indicated nonnegative domains, so
the nonlinear expression does too. Apply the three upper inequalities
to $F_\theta$, then substitute this proved initial upper bound.
This proves the propagated budget at the original energy use site:
\[
 2\sum_{N=q-1}^{2q-1}w_N
 \log\max\left\{1,\frac{\mathcal L}{\sqrt{2J_N^D}}-1\right\}
 \leq\mathcal B_{k,\theta}\leq U_{k,\theta;p,r}.
\tag{HC13a}
\]
The $w_N$, $J_N^D$ and every original derivative-source and boundary
term in (HC12)--(HC13) remain unchanged. Coefficient zero insertion
at a specified $(W,S;A)$ commutes with the displayed inclusions and
linear source maps. The minimum-section comparison retains the
original monic-division primitive and the proper-source residual
proved in (ISM38)--(ISM43); no residual class or empty mask's outer
label is set to zero by this metric calculation.
For the averaged original phase metrics the already proved stronger
comparison $(1-\eta^2)G_N\preceq\overline G_N\preceq G_N$ gives instead
$|\overline{\mathcal B}_k-\mathcal B_k|
\leq2q\log(1/(1-\eta^2))$.  These changes are of order $q$ for a fixed
$\eta$, whereas the compatible width (HC11) has positive quadratic
order.  The exact Laplacian residual (HT17) remains present at every
phase.  No displayed upper estimate for its norm, or for the original
derivative moments in (HC12), contradicts these lower bounds.

\subsection{The sharper finite comparison retains its exact error}

The original lower-envelope measure is
\[
 d\nu_k(u)=c_h\vartheta_h^{k-3}e^{-\alpha(k-3)}
 e^{-\alpha|u|}(k-2+|u|)^{-B}\,du.
\]
Let $V_N^\nu$ be its quotient volume for the identical full remainder
map.  The actual finite upper expression for the four-window volume is
\[
 \mathcal U_{\nu,k}=\log(V_{q-1}V_q)
                  -\log(V_{2q-1}^\nu V_{2q}^\nu),
\]
with the exact nonnegative error
\[
 \boxed{
 \mathcal U_{\nu,k}-\mathcal B_k
 =\log\frac{V_{2q-1}}{V_{2q-1}^\nu}
  +\log\frac{V_{2q}}{V_{2q}^\nu}\geq0.}\tag{HC15}
\]
Indeed $m_{h,k}(u)\,du\geq d\nu_k$ bounds every source polynomial
quadratic form.  Minimizing on each identical affine remainder fibre
gives $G_N\succeq G_N^\nu\succ0$.  Conjugation by
$(G_N^\nu)^{-1/2}$ shows that each volume ratio in (HC15) is at least
one.  Cancellation proves the equality.  The low-degree factors require
only original moments through $2q$, and the comparison moments in the
last factors go through $4q$.  These are well-defined finite expressions,
with the actual full jet map and source mass.  The supplied proof
records do not evaluate their errors at the hypothetical arithmetic
quartet or bound them at a rate opposing (HC5).  Retaining this sharper
finite expression therefore supplies no extra contradictory inequality.
